\section{Tracial approximation on a cone}\label{tc:section}

Throughout this section, $Y=(0,1]$, $m$ is a finite Borel measure of full
support on $Y$, and $S=C_0(Y)\otimes P$.  We retain the invariant maps
$F:T(Q)\to T(P)$ and $G:\Cu(P)\to\Cu(Q)$, so that
\begin{equation}\label{tc:transport}
 d_\tau(G[a])=d_{F(\tau)}(a).
\end{equation}
We seek a completely positive map that is multiplicative modulo the
local error ideals.  It must carry two kinds of information.  Its square
moments will identify the prescribed trace wherever that trace is densely
finite.  The ideal generated by its image of a nonzero elementary tensor
must also contain the corresponding constant ideal.  This second condition will prevent the
corrected map from turning an infinite part of a weight into a finite one.

Matrix traces are unnormalized.  The rational averages used below have a
separate, normalized UHF tensor factor; \Cref{st:folding} identifies their
ranges with $Q$ with precisely this trace convention.

Choose a countable family $\mathcal E$ of positive elementary tensors
$s=f\otimes a$ as follows.  In each ideal in the countable basis, and in
each finite union of those ideals, take the positive rational cuts of a
norm-dense set of positives.  Retain slightly larger positive pre-cuts in the
same ideal.  Do the same for positive functions in $C_0(Y)$, and include
the powers and further cuts of the factors that occur below.  Thus the
local linear spans of $\mathcal E$ are dense, and its rectangle supports
exhaust the open subsets of $Y\times X$.  All elementary tests used in
this section may be included in this one countable family.

For a test $s_i=f_i\otimes a_i$ supported over an open set $U_i$, choose,
once and for all, a positive controller $b_i\in Q(U_i)$ of compact Cuntz
support such that
\begin{equation}\label{tc:controller}
 G[a_i]\leq[b_i].
\end{equation}
The larger pre-cuts allow this choice inside $Q(U_i)$, with room for a
further compactly supported upper controller there.  When needed, the
same test is listed with different containing basis opens.  This ensures
that the controllers required for a local assertion lie in that open,
as in \Cref{seq:diagonal}.

\begin{theorem}\label{tc:rough-model}
There is a completely positive contraction $h:S\to D=Q_\infty$ with the
following properties.

For every open $U\subseteq X$, $s\in S(U)$, and $z\in S$,
\begin{equation}\label{tc:mixed}
 h(sz)-h(s)h(z),\quad h(zs)-h(z)h(s)\ \in J(U).
\end{equation}
If $s=f\otimes a\in\mathcal E$ is controlled inside $U$, and $\rho$ is a
regularized limit trace whose restriction $\sigma=\rho|_Q$ is densely
finite on $Q(U)$, then
\begin{equation}\label{tc:square}
 \rho(h(s)^2)=\left(\int_Y f^2\,dm\right)F(\sigma)(a^2).
\end{equation}
Finally, for every $f\otimes a\in\mathcal E$ with $f\ne0$,
\begin{equation}\label{tc:lower-ideal}
 Q(U_a)\subseteq\Ideal_D\bigl(h(f\otimes a)\bigr).
\end{equation}
The test family may be enlarged in advance by any countable family of
the controlled elementary tests just described.
\end{theorem}

For a finite-stage completely positive contraction $H:S\to M_k(Q)$,
we keep three trace coordinates for each positive input $s$: the trace
values of $H(s)$, $H(s^2)$, and $H(s)^2$.  Matching the last
two makes the positive Schwarz defect small in trace.  The first bounds
the square mass that could otherwise disappear when a coordinate trace
is regularized.

Trace approximation alone does not give the lower ideal inclusion.  For
that we also approximate positive expressions
$yv^*g(H(s))vy$ by $y^2$, where $y$ is a positive constant from the prescribed
ideal of $Q$, $v\in M_{k,1}(Q)$ is a contraction, and the continuous function
$g:[0,\infty)\to[0,1]$ vanishes below a fixed positive cutoff.  Their norm approximation will force a lower rank bound even at
weights for which the controller has infinite rank.  We first realize
these finite data in von Neumann algebras, then approximate them in $Q$
and average.  A final diagonal choice retains the fixed controllers and
cutoffs.

\subsection{Realization in von Neumann algebras}

We first record explicitly the trace-preserving embedding needed for
extended weights. The embedding is not required to be unital.
We write $R$ for the hyperfinite factor of type $\mathrm{II}_1$ and
$R_\infty=R\,\overline\otimes\,B(\ell^2)$ for the hyperfinite factor
of type $\mathrm{II}_\infty$. Their uniqueness is part of Connes's
classification \cite{Connes}; see also the distinct contemporary
report \cite[Theorems~4.2.1 and~4.3.1]{ConnesReport}.

\begin{lemma}[Trace-preserving semifinite embedding]\label{tc:semifinite-embedding}
Let $N$ be a hyperfinite von Neumann algebra with separable predual and a
faithful normal semifinite trace $T$.  There is a normal injective
homomorphism from $N$ into the hyperfinite factor $R_\infty$ of type
$\mathrm{II}_\infty$ preserving $T$.  Consequently, such an embedding
exists into any prescribed infinite corner of a separably acting factor
of type $\mathrm{II}_\infty$.
\end{lemma}

\begin{proof}
The central decomposition of a semifinite hyperfinite algebra gives a
countable direct sum of algebras
\[
 L^\infty(X_j,\mu_j^0)\,\overline\otimes\,N_j,
 \qquad
 N_j\in\{M_n,\ B(\ell^2),\ R,\ R_\infty\}.
\]
Here the measure spaces are standard and the factor traces are the
usual matrix trace or the standard finite or semifinite trace.  This is
the semifinite form of the hyperfinite central decomposition; see
\cite[Section~6, p.~49, after (6.48) and (6.50)]{HRS}.  Relative to these
traces and a finite measure $\mu_j^0$ in the measure class, $T$ has a
central Radon--Nikodym density $h_j$.  Faithfulness and semifiniteness
give $0<h_j<\infty$ almost everywhere.  Replacing $\mu_j^0$ by
$\mu_j=h_j\mu_j^0$ absorbs the density and leaves a standard
$\sigma$-finite measure space.

There is a normal, trace-preserving injection
\[
 L^\infty(X_j,\mu_j)\longrightarrow
 L^\infty([0,\infty),dt).
\]
Indeed, partition $X_j$ into countably many pieces of finite measure.
On each piece assign an interval of the appropriate length to each
atom; the remaining diffuse part is isomorphic, modulo null sets, to
the remaining interval.  Pullback by the resulting measure-preserving
map, extended by zero outside the union of the intervals, gives the
injection.  Moreover,
$L^\infty([0,\infty))$ embeds trace-preservingly as
$L^\infty([0,1])\overline\otimes\ell^\infty$ in
$R\overline\otimes B(\ell^2)=R_\infty$.

Each standard factor $N_j$ embeds trace-preservingly in $R_\infty$:
use $1_R\otimes M_n$, $1_R\otimes B(\ell^2)$,
$R\otimes e_{11}$, or the identity map, respectively.  Tensor these
embeddings with the preceding abelian embeddings and use the
trace-preserving isomorphism
$R_\infty\overline\otimes R_\infty\cong R_\infty$.
The countably many summands can be placed in the diagonal corners of
$R_\infty\overline\otimes B(\ell^2)$.  The orthogonal strong sum is
normal, injective, and trace preserving.

Finally, in an infinite corner of an arbitrary factor of type
$\mathrm{II}_\infty$, take a projection of trace one, nested unital
matrix algebras $M_{2^n}$ in its finite corner, and their weak closure.
They give a trace-preserving copy of $R$.  Matrix amplification in the
infinite corner gives $R_\infty$.  Composition completes the proof.
\end{proof}

\begin{lemma}[Realization in a factor]\label{tc:factor-realization}
Let $\pi:Q\to M$ be a nondegenerate representation with weakly dense
range in a separably acting factor.  There is a homomorphism
$l:S\to M$ with the following rank data.  If $M$ is semifinite with
faithful dimension trace $T$, then
\begin{equation}\label{tc:factor-rank}
 d_T(l(f\otimes a))
   =m(\{f>0\})d_{F(T\circ\pi)}(a)
 \qquad(f,a\geq0).
\end{equation}
In a type III factor the corresponding support projection is nonzero
exactly when the transported ideal data prescribe a nonzero support.
\end{lemma}

\begin{proof}
Stability of $Q$ supplies multiplier isometries whose normal images
show that $M$ is properly infinite. Since $M$ is separably acting,
the usual scalar trace-dimension comparison applies in its semifinite
factor cases; in its type~III case all nonzero projections are
equivalent \cite[III.1.7.9--10 and III.1.3.7]{Blackadar}.  Suppose first that $M$ is of type
I.  The rank function $d_{T\circ\pi}$ on $\Cu(Q)$ takes only the values
$0$ and $\infty$.  Otherwise real scalar divisibility of $\Cu(Q)$,
applied to a class of finite nonzero integer rank, would give a class
of noninteger rank in a type I factor.  Thus the right side of
\eqref{tc:factor-rank} also takes only these two values.  Its zero ideal
on $P$ is the ideal transported from $\ker\pi$.  A faithful separable
representation of the corresponding quotient of $S$, repeated with
infinite multiplicity and embedded in $M$, realizes the required ranks.
For type III the same construction realizes the zero ideal; all
nonzero support projections are equivalent.  The assertion is also
immediate when the prescribed quotient is zero.

Suppose now that $M$ has type $\mathrm{II}_\infty$, and put
\[
 \nu=m\otimes F(T\circ\pi),\qquad I=\Fin(\nu),
 \qquad K=\Ker(\nu).
\]
Here $\Fin(\nu)$ is the closed ideal on which $\nu$ is densely finite,
and $K\subseteq I$ is its zero ideal.  The trace representation of
$I/K$ has weak closure $N$ carrying the faithful normal semifinite
trace induced by $\nu$.  Nuclearity of $I/K$ makes $N$ hyperfinite
\cite{Connes}; the nuclear-representation consequence is also stated
in \cite[Theorem~4.1 and the following paragraph]{ConnesReport}.
The algebra $N$ has separable predual.  Its nondegenerate
representation extends from $I$ to $S$ by the action of $S$ on the
ideal $I$.  \Cref{tc:semifinite-embedding} embeds this representation
trace-preservingly in one infinite corner of $M$.

In an orthogonal infinite corner, choose a copy of $B(\ell^2)$ whose
minimal projections have a fixed positive finite trace.  In this copy
place an infinite multiple of a faithful separable representation of
$S/I$.  On $I$ this second summand is zero.  If $x\in S_+$ does not
belong to $I$, its second image has infinite trace and infinite support
rank.  This is exactly what is required, since $\nu(x)=\infty$ off
$I$.  On $I$, the first summand gives $\nu$ and all its ranks, including
infinite values.  The combined homomorphism therefore satisfies
$T\circ l=\nu$, has zero ideal $K$, and proves
\eqref{tc:factor-rank}.  If $I=0$ or $I=S$, simply omit the empty
summand.
\end{proof}

Factorwise existence is enough only if the maps can be assembled
measurably.  We encode their rank requirements by partial isometries
between support projections.  These witnesses assemble along with the
maps, and later allow each selected normal trace to evaluate the required
ranks directly.

\begin{lemma}[Measurable realization]\label{tc:measurable-realization}
Let $\pi$ be a nondegenerate separably acting representation of $Q$,
and put $M=\pi(Q)''$.  There is a homomorphism $l:S\to M$ whose
factorwise ranks are those in \Cref{tc:factor-realization}, on any
prescribed countable elementary cut-test family.  In particular, it
may be chosen on the entire family $\mathcal E$ and its required cuts.
\end{lemma}

\begin{proof}
Use the standard central decomposition of $M$, modulo null sets, into
separably acting factors $M_x$ \cite[III.5.1.14--15]{Blackadar}.  We explain the measurable selection
in a form that does not choose a normalization of the factor traces.
For each elementary test $s_j=f_j\otimes a_j$ fix $c_j\in Q_+$ with
\begin{equation}\label{tc:transported-representative}
 [c_j]=m(\{f_j>0\})G[a_j].
\end{equation}
Realification and stability provide these representatives.  Code a
contractive homomorphism on a countable dense rational $*$-algebra in
$S$ by its values in the strong-$*$ bounded balls of the factor field.
Include norm-convergent approximations to the chosen positive tests.
Multiplication on bounded balls, adjoints, norm bounds, and membership
in the measurable algebra field give countably many Borel conditions.
Membership may be checked against the strong-$*$ dense measurable
sections of the field.  Adjoin contraction coordinates $w_j\in M_x$
subject to
\begin{equation}\label{tc:selection-witness}
 w_j^*w_j=\supp(l_x(s_j)),\qquad
 w_jw_j^*=\supp(\pi_x(c_j)).
\end{equation}
The support map is the strong limit of
$t\mapsto t(t+1/n)^{-1}$, so these are again Borel conditions.

Thus the augmented admissible-code relation is Borel in a countable
product of standard measurable bounded operator fields. Its sections
are nonempty by \Cref{tc:factor-realization}: equal factor dimensions
give the partial isometries in \eqref{tc:selection-witness}. Apply
Jankov--von Neumann measurable uniformization to this augmented analytic
relation, retaining both the homomorphism coordinates and
all the contraction witnesses $w_j$. The selector is measurable for
the completed central measure. Replace its countably many coordinates
by measurable representatives and discard one common null set on which
any selected algebraic equation or support identity fails.
This is precisely the measured-field setting, as distinguished from
an everywhere Borel field in \cite[Section~8, especially Lemma~8.6(1),
and the proof of Proposition~3.2]{VaesWouters}.
The uniform operator bounds make the selected values, including the
$w_j$, decomposable operators. Their countably many algebraic relations
define a homomorphism $l:S\to M$, and the assembled witnesses establish
\eqref{tc:selection-witness} in $M$. Every elementary cut test needed
below is included in the code, so its prescribed rank data hold in
this field.
\end{proof}

Here is the finite-trace consequence of this construction that we will
use.  It specifies the normal trace associated to each selected weight
and the finite constants through which the rank data are evaluated.

\begin{lemma}[Realization for finitely many weights]\label{tc:finite-trace-realization}
Fix finitely many controlled tests $s_i=f_i\otimes a_i$ and finitely
many traces $\tau$ for each $i$, each satisfying $d_\tau(b_i)<\infty$.
Take a separably acting representation $\pi$ containing the induced
trace representations of their densely finite ideals, as well as any
prescribed finite family of representations for bounded functionals
on $Q$.  The homomorphism $l$ in \Cref{tc:measurable-realization} may
be chosen so that, for the induced normal trace $\widetilde\tau$ on
$M=\pi(Q)''$,
\begin{equation}\label{tc:finite-values}
 \widetilde\tau(l(s_i^k))
   =\left(\int f_i^k\,dm\right)F(\tau)(a_i^k),\qquad k=1,2.
\end{equation}
This assertion requires no commutation of $F$ with a direct integral
of extended weights.
\end{lemma}

\begin{proof}
Fix one selected $\tau$ and put $I=\Fin(\tau)$.  Its trace representation
$\pi_\tau$ on $I/\Ker(\tau)$ has weak closure $N_\tau$ with faithful
normal semifinite trace $T_\tau$.  The nondegenerate extension of
$\pi_\tau$ to $Q$ still has image in $N_\tau$: for an approximate unit
$(u_\lambda)$ of $I$, the operators
$\pi_\tau(q)\pi_\tau(u_\lambda)\in\pi_\tau(I)$ converge strongly to
$\pi_\tau(q)$.  Hence $\pi_\tau(Q)''=N_\tau$.

The projection onto this representation summand commutes with
$M=\pi(Q)''$.  Restriction to that summand is therefore a normal
surjective homomorphism $\theta_\tau:M\to N_\tau$.  Its kernel has
the form $(1-z_\tau)M$ for a central projection $z_\tau$, and its
restriction to $z_\tau M$ is a normal isomorphism.  Define
\[
 \widetilde\tau=T_\tau\circ\theta_\tau.
\]
This is a normal semifinite trace, zero on $(1-z_\tau)M$, and it agrees
with $\tau$ on $\pi(I)_+$.  Normality also gives
$d_{\widetilde\tau}(\pi(c))=d_\tau(c)$ for $c\in I_+$.
The projections $z_\tau$ for different selected weights may overlap:
each trace is defined through its own normal quotient, with no
disjointness or additional contribution from other representation
summands.  The construction includes the zero trace representation.
Agreement with $\tau$ outside $I$ is neither asserted nor needed.

Include the rectangle tests
$(f_i-r)_+\otimes(a_i-t)_+$ for rational $r,t\geq0$, and choose
$c_{r,t}\in Q_+$ with
\[
 [c_{r,t}]=m(\{f_i>r\})G[(a_i-t)_+].
\]
Their ranks are at most $m(Y)d_\tau(b_i)<\infty$, so $c_{r,t}\in I$.
Retaining the selected witness coordinates in
\eqref{tc:selection-witness} gives, in $M$, equivalence of the support
of $l((f_i-r)_+\otimes(a_i-t)_+)$ with $\supp(\pi(c_{r,t}))$.
Applying $\widetilde\tau$ directly to these equivalent projections gives
\[
 \widetilde\tau\bigl(\supp l((f_i-r)_+\otimes(a_i-t)_+)\bigr)
    =m(\{f_i>r\})d_{F(\tau)}((a_i-t)_+).
\]
In particular, $\supp(l(s_i))$ has finite $\widetilde\tau$-trace.

In the support corner of $l$, the multiplier images of $f_i\otimes1$
and $1\otimes a_i$ commute, and their product is $l(s_i)$.  The last
display specifies their finite joint spectral measure on every upper
rectangle in the positive quadrant.  It is the product of the
spectral measure of $f_i$ under $m$ and that of $a_i$ under $F(\tau)$.
Increasing positive step approximations to the product function
$(u,v)\mapsto u^kv^k$, using rational spectral partitions away from
zero, therefore give \eqref{tc:finite-values} by normal monotone
convergence, for $k=1,2$.  All rank identities used here were evaluated
on constants in $\Fin(\tau)$.
\end{proof}

\subsection{Simultaneous finite approximation}

The von Neumann algebra realizations attain the prescribed trace and
support data. We now approximate finitely many of those data inside
$Q$. The approximations must retain a finite rank bound over an entire
trace cap, since pointwise trace approximation would not control the
local error ideals after diagonalization.

Fix finitely many tests $s_i=f_i\otimes a_i$, indexed by $1\leq i\leq n$,
with their fixed controllers.  Set
\begin{equation}\label{tc:caps-targets}
 K_i=\{\tau\in T(Q):d_\tau(b_i)\leq1\},\qquad
 t_i^{(k)}(\tau)=\left(\int f_i^k\,dm\right)F(\tau)(a_i^k).
\end{equation}
By \Cref{un:trace-cap}, $K_i$ is compact and convex.  The functions
$t_i^{(k)}$ are finite continuous affine functions on $K_i$.

We impose finitely many lower comparisons, whose indices $j$ refer to
some of the same tests.  Choose $\delta_j>0$ and $r_j>0$ with
\begin{equation}\label{tc:lower-data}
 r_j<m(\{f_j>\delta_j\}),\qquad
 [y_j]=r_jG[(a_j-\delta_j)_+],\qquad y_j\in Q_+.
\end{equation}
Rescale representatives so that $\|y_j\|\leq1$.  Choose a continuous
function $g_j:[0,\infty)\to[0,1]$ and constants
\begin{equation}\label{tc:fixed-slack}
 0<\eta_j<\theta_j<\delta_j^2,
 \qquad g_j|_{[0,\eta_j]}=0,
 \qquad g_j|_{[\theta_j,\infty)}=1.
\end{equation}
These constants are fixed before any finite approximation or averaging.

A candidate consists of a completely positive contraction
$H:S\to M_k(Q)$ and contractions $v_j\in M_{k,1}(Q)$, with the following
whole-cap support condition: for each $i$ there is a finite constant
$C_{H,i}$ such that, for every $\tau\in T(Q)$,
\begin{equation}\label{tc:candidate-rank}
 \begin{split}
 d_{\tau\otimes\Tr}(H(s_i))&\leq C_{H,i}d_\tau(b_i),\\
 d_{\tau\otimes\Tr}(H(s_i^2))&\leq C_{H,i}d_\tau(b_i).
 \end{split}
\end{equation}
The constants may depend on the candidate.  Its vector in the real
locally convex space
\[
 \prod_{i=1}^n C(K_i,\R)^3\ \times\ \prod_j Q_{\mathrm{sa}}
\]
is
\begin{equation}\label{tc:candidate-vector}
 \left(
 (\tau(H(s_i)),\tau(H(s_i^2)),\tau(H(s_i)^2))_i;
 (y_jv_j^*g_j(H(s_j))v_jy_j)_j
 \right).
\end{equation}
Here and below $\tau$ on matrices means $\tau\otimes\Tr$.
The support bounds and \Cref{un:trace-cap} make all three trace
coordinates continuous.  No bound uniform over all candidates is
needed for this assertion.

\begin{proposition}[Simultaneous finite approximation]\label{tc:finite-convexity}
The target vector
\begin{equation}\label{tc:target-vector}
 \bigl((t_i^{(1)},t_i^{(2)},t_i^{(2)})_i;(y_j^2)_j\bigr)
\end{equation}
belongs to the closed convex hull of the candidate vectors.
\end{proposition}

\begin{proof}
\smallskip\noindent\emph{Reduce a separator to finitely many weights.} 
Suppose a continuous real linear functional strictly separates the
target from this convex hull.  On each $C(K_i,\R)$ coordinate it is
integration against a finite signed measure.  Split each measure into
its positive and negative parts.  By \Cref{un:barycenter}, each nonzero
part has a barycenter in $K_i$ reproducing every finite continuous
affine evaluation used here; its mass is an ordinary scalar
coefficient.  This holds simultaneously for the candidate evaluations
and the target evaluations, including candidates with different
support constants.  Thus the separating functional involves only
finitely many selected traces.  On the algebra coordinates it involves
finitely many bounded selfadjoint functionals.

\smallskip\noindent\emph{Realize the selected values exactly.} 
Choose a faithful separably acting representation $\pi$ of $Q$
containing representations for these bounded functionals and the
induced representations of the densely finite ideals of the selected
traces.  Write $M=\pi(Q)''$ and identify $Q$ with its image.
\Cref{tc:measurable-realization,tc:finite-trace-realization} give a
homomorphism $l:S\to M$ attaining the selected trace values.  We next
attain the lower norm coordinates in $M$.

Let $q_j$ be the support projection of
$l((f_j-\delta_j)_+\otimes(a_j-\delta_j)_+)$.  Functional calculus in
the two commuting factors gives
\[
 q_j\leq1_{(\delta_j^2,\infty)}(l(s_j)).
\]
The rank data and \eqref{tc:lower-data} give, in every factor,
$\supp(y_j)\precsim q_j$.  This comparison also covers zero and
infinite dimensions and the type III parts.  Comparison of projections
in a direct integral therefore gives a partial isometry $v_j\in M$
with initial projection $\supp(y_j)$ and final projection at most
$q_j$.  Hence
\begin{equation}\label{tc:exact-norm-signal}
 y_jv_j^*g_j(l(s_j))v_jy_j=y_j^2.
\end{equation}

\smallskip\noindent\emph{Approximate inside fixed supports.} 
It remains to approximate these finite data by candidates satisfying
\eqref{tc:candidate-rank} on the entire caps.  A tested value must vanish
exactly in a block where its controller vanishes; strong approximation
alone would not retain this condition.  We therefore split into central
sectors and factor each sector map through the required source quotient
before approximating it.  Let $z_i$ be the central
support of $b_i$ in $M$, and form the finitely many central sectors
\[
 z_E=\prod_{i\in E}z_i\prod_{i\notin E}(1-z_i),
 \qquad E\subseteq\{1,\ldots,n\}.
\]
If $i\notin E$, the rank data and $G[a_i]\leq[b_i]$ imply
$z_El(s_i)=0$.  Let $J_E^S$ be the ideal of $S$ generated by these
$s_i$, and let
\[
 I_E=\bigcap_{i\in E}\Ideal_Q(b_i),
 \qquad I_{\varnothing}=Q.
\]
The weak closure of $I_E$ has central support dominating $z_E$.
Choose positive contractions $e_E\in I_E$ with compact Cuntz support,
tending strongly to the identity on this sector.  The maps
\begin{equation}\label{tc:sector-compression}
 \psi_E(s)=e_Ez_El(s)e_E
\end{equation}
take values in $\Her(e_E)''$, are completely positive contractions,
and factor through $S/J_E^S$.  This quotient is taken before the
approximations below.  The element $e_E$ itself need not vanish outside
the sector; \eqref{tc:sector-compression} does vanish there.
Write
\[
 D_0(s)=\operatorname{diag}_E(z_El(s)),\qquad
 D_e(s)=\operatorname{diag}_E(e_Ez_El(s)e_E).
\]
Thus $D_0$ is the exact sector representation and $D_e$ is its completely
positive compression.  As the $e_E$ increase along approximate units,
$D_e(s)\to D_0(s)$ strongly-$*$ for each $s$.

For the selected trace belonging to test $i$, the sector contributes
zero if $i\notin E$.  If $i\in E$, the positive operators
$z_El(s_i)$ and $z_El(s_i^2)$ are integrable, and $z_El(s_i)$ is square
integrable, by \eqref{tc:finite-values}.  Strongly converging bounded
approximate units therefore make the compressions converge in the
required $L^1$ and $L^2$ seminorms.  For the $L^1$ assertion apply
$L^2$ convergence to the square root of the integrable positive
operator and use Cauchy--Schwarz; for the square evaluation apply it
directly to $z_El(s_i)$.

For the lower signal set $V_j=(z_Ev_j)_E$.  This column is a
contraction, and, because $g_j(0)=0$, the exact sector representation
satisfies
\[
 y_jV_j^*g_j(D_0(s_j))V_jy_j=y_j^2.
\]
Bounded strong-$*$ functional calculus gives
$g_j(D_e(s_j))\to g_j(D_0(s_j))$.  Fix the $e_E$ sufficiently far out
that replacing $D_0$ by $D_e$ makes the errors small against the selected
normal functionals, as well as meeting the preceding finite trace
requirements.  The compressed expression need not equal $y_j^2$ exactly.

We give the remaining approximation explicitly. Nuclearity of the
quotient $R_0=S/J_E^S$ provides completely positive contractive
factorizations through finite matrix algebras approximating its
identity on the finite tests \cite{ChoiEffros}; see also
\cite[Corollary~3.3]{HanPaulsen}. For nonunital $R_0$, apply the latter
unital statement to its nuclear unitization. Restrict the incoming
map to $R_0$ and compress the outgoing map $\beta$ to
$x\mapsto e\beta(x)e$, where $e\in R_0$ is a positive contractive
approximate-unit element. Choose $e$ first so that
$\|eae-a\|<\delta/2$ on the finite tests, and then choose the
factorization with error less than $\delta/2$ there. Contractivity
bounds the total error by $\delta$, and the compressed map takes
values in $R_0$. Compose this outgoing map with $\psi_E$.
A completely positive map from $M_d$ into $\Her(e_E)''$ is specified
by its positive Choi matrix.  Approximate the square root of that
matrix strongly-$*$, with bounded norm, by matrices over
$\Her(e_E)$; the products give positive Choi matrices over the
hereditary algebra.  If the resulting map $\beta$ has
$t=\beta(1)$, replace it by
\[
 x\longmapsto \max(1,t)^{-1/2}\beta(x)\max(1,t)^{-1/2}.
\]
The functional calculus is in the unitization.  This replacement is
completely positive and contractive, stays in $\Her(e_E)$, and
retains the strong-$*$ limit, because the limiting unit value is at
most one.  We have obtained completely positive contractions
\[
 H_E:S/J_E^S\longrightarrow\Her(e_E)\subseteq Q
\]
approximating \eqref{tc:sector-compression} on the required tests.

The support element $e_E$ was fixed before this Choi-matrix
approximation, and every approximating value remains in $\Her(e_E)$.
For each $i\in E$, compact Cuntz support inside $I_E$ gives a finite
integer $N_{E,i}$ with $[e_E]\leq N_{E,i}[b_i]$. Thus the support of
every output in this block has rank at most
$N_{E,i}d_\tau(b_i)$, for every extended trace $\tau$.  On the
selected traces this is one fixed finite normal trace support.
Bounded strong-$*$ convergence on that support gives convergence of
the trace and square-trace evaluations, with no loss of mass. For
$i\notin E$, the quotient-first factorization kills $s_i$ and $s_i^2$
exactly, so no trace-convergence argument is required in that block.
This proves the trace approximation, as well as the required
uniformity on the whole caps.

Put the sector maps in separate matrix blocks:
$H=\bigoplus_E H_E$.  This is a completely positive contraction since
the norm of a direct sum is the maximum of the norms.  Its
unnormalized matrix trace is the sum of the sector traces.  The blocks
with $i\notin E$ kill $s_i$ and $s_i^2$ exactly, while
\[
 d_\tau(H(s_i)),\ d_\tau(H(s_i^2))
   \leq\left(\sum_{E\ni i}N_{E,i}\right)d_\tau(b_i).
\]
The maps $H_E$ approximate the fixed compressed map $D_e$,
not the exact map $D_0$.  Rectangular Kaplansky density approximates the
columns $V_j$ strongly-$*$ by contraction columns over $Q$.  Bounded
strong-$*$ functional calculus and multiplication therefore make the
resulting algebra coordinates arbitrarily close, against the selected
normal functionals, to those for $D_e$.  Combined with the preceding
choice of the $e_E$, this approximates $y_j^2$ to the required functional
accuracy.  Those functionals are normal in $\pi$ by construction.

Thus candidate vectors approximate the value of every finite
functional obtained from a supposed separator, arbitrarily closely
to the target value.  This contradicts strict separation and proves
the proposition.
\end{proof}

\subsection{Averaging and the lower rank bound}

The preceding proposition gives finite convex combinations with
arbitrarily small uniform errors on the caps and arbitrarily small
norm errors in the algebra coordinates.  The convex coefficients can
be made rational: only finitely many bounded vectors are involved in
any one combination.  If their common denominator is $N$, put the
candidates, with their rational multiplicities, in the $N$ diagonal
blocks of a normalized matrix subalgebra of $\UHF$.  First zero-pad
their internal matrix spaces to a common size.  Denote the resulting
completely positive contraction by $H_{\mathrm{av}}$.  Then, for
$k=1,2$,
\begin{equation}\label{tc:average-traces}
 \begin{split}
 \tau(H_{\mathrm{av}}(s_i^k))
    &=\sum_\alpha\lambda_\alpha\tau(H_\alpha(s_i^k)),\\
 \tau(H_{\mathrm{av}}(s_i)^2)
    &=\sum_\alpha\lambda_\alpha\tau(H_\alpha(s_i)^2).
 \end{split}
\end{equation}
Only the $\alpha$ direction is normalized.  Internal sector and
stabilizing matrix directions keep the unnormalized trace.  Folding
as in \Cref{st:folding} therefore gives the same formulas for a map
with range $Q$.

\begin{lemma}[A lower rank bound preserved by averaging]\label{tc:rank-average}
Fix lower data $y,g,\eta$ as in
\eqref{tc:lower-data}--\eqref{tc:fixed-slack}.  If a finite rational
average satisfies
\[
 \left\|\sum_\alpha\lambda_\alpha
       yv_\alpha^*g(H_\alpha(s))v_\alpha y-y^2\right\|
       <\eps^2/2,
\]
then for every extended trace on $Q$,
\begin{equation}\label{tc:rank-bound}
 d_\tau((y-\eps)_+)
   \leq \frac{2\|y\|^2}{\eps^2}\,
        d_\tau((H_{\mathrm{av}}(s)-\eta)_+).
\end{equation}
The constant and the cut $\eta$ are independent of the number of
candidates, their multiplicities, and their internal matrix sizes.
\end{lemma}

\begin{proof}
Write the displayed average as $T$ and let
$p=1_{(\eps,\infty)}(y)$ in the bidual.  The norm estimate gives
$pTp\geq(\eps^2/2)p$.  Since the $v_\alpha$ are contractions and
$0\leq g(t)\leq1_{(\eta,\infty)}(t)$, positivity and traciality give
\[
 \frac{\eps^2}{2}\tau(p)
 \leq \|y\|^2\sum_\alpha\lambda_\alpha
       d_\tau((H_\alpha(s)-\eta)_+).
\]
The sum is exactly the rank on the right of
\eqref{tc:rank-bound}, by the normalized averaging convention.

Here is a justification when $\tau$ is not densely finite.  If the
right side is infinite, there is nothing to prove.  If it is finite,
discard zero coefficients.  Every support of $g(H_\alpha(s))$ then
has finite rank. The norm estimate gives
$(y^2-\eps^2/2)_+\precsim T$, and $T$ is Cuntz below the finite
direct sum of the $g(H_\alpha(s))$, by its finite row factorization.
Since $p$ is the support of $(y^2-\eps^2)_+$, it has finite rank
and belongs to the normal trace's finite part. One may initially
use a slightly smaller spectral projection inside $p$.
The displayed computation is now a computation for the induced
normal trace on this finite part.  Let the additional cut slack decrease to zero.  This gives
the inequality for $p$ by normal monotone convergence and proves
the assertion for every extended weight.
\end{proof}

\subsection{The diagonal choice}

\begin{proof}[Proof of \Cref{tc:rough-model}]
\smallskip\noindent\emph{Fix the comparisons before forming averages.} 
Fix the test family, its larger pre-cuts, its local controllers, and
all lower data before choosing the sequence.  For each nonzero
$f\otimes a\in\mathcal E$, take countably many pairs $(\delta,r)$
with
\[
 \delta\downarrow0,\qquad 0<r<m(\{f>\delta\}),
\]
and positives $y$ representing $rG[(a-\delta)_+]$.  Include all
positive rational cuts of these $y$ in the lower requirements.  Their
ideals exhaust $Q(U_a)$: multiplication by a positive real scalar
does not change a Cuntz ideal, the cuts of $a$ exhaust its ideal, and
$G$ identifies ideals.  For each requirement fix its function $g$,
its lower cutoff $\eta$, and its source cutoff $\eps$.  In particular,
none of these cutoffs depends on the eventual averaging denominator.

At stage $n$, apply \Cref{tc:finite-convexity} to the first $n$ trace
tests and first $n$ lower requirements.  Choose rational averages and
fold them into $Q$, obtaining a completely positive contraction
$h_n:S\to Q$.  Given any prescribed sequence $\gamma_n\downarrow0$,
the trace errors can be made to satisfy, for $i\leq n$,
\begin{equation}\label{tc:coordinate-errors}
 \begin{split}
 |\tau(h_n(s_i))-t_i^{(1)}(\tau)|
        &\leq\gamma_n d_\tau(b_i),\\
 |\tau(h_n(s_i^2))-t_i^{(2)}(\tau)|
        &\leq\gamma_n d_\tau(b_i),\\
 |\tau(h_n(s_i)^2)-t_i^{(2)}(\tau)|
        &\leq\gamma_n d_\tau(b_i)
 \end{split}
\end{equation}
whenever $d_\tau(b_i)<\infty$.  Indeed, the uniform estimate on $K_i$
rescales when this rank is nonzero; when it is zero, apply the cap
estimate to every positive scalar multiple of $\tau$ and let the
scalar increase.  All evaluations in that argument are finite, by
the support bounds.  Simultaneously choose the algebra norm errors
below $\eps^2/2$ for each of the first $n$ fixed lower requirements.
\Cref{tc:rank-average} then supplies the corresponding lower rank
inequalities with their fixed constants.

\smallskip\noindent\emph{Localize the multiplicative defects.} 
The coordinate maps define a completely positive contraction
$h:S\to D$.  Put
$\Delta_{n,i}=h_n(s_i^2)-h_n(s_i)^2\geq0$ by Schwarz's inequality.
For $t>0$, \eqref{tc:coordinate-errors} gives
\begin{equation}\label{tc:defect-rank}
 d_\tau((\Delta_{n,i}-t)_+)
       \leq\frac{2\gamma_n}{t}d_\tau(b_i).
\end{equation}
For zero controller rank this follows by the same rescaling argument;
for infinite controller rank the upper bound is automatic.  Thus
$h(s_i^2)-h(s_i)^2\in J(U_i)$.  In $D/J(U_i)$ the positive element
$s_i$ belongs to the multiplicative domain of the induced completely
positive contraction.  Hence both mixed defects with $s_i$ belong
to $J(U_i)$, for every $z\in S$.

For an arbitrary open $U$, use the tests localized in basis ideals
inside $U$.  Their linear spans are dense in $S(U)$.  Inclusion and
closedness of the ideals $J(V)\subseteq J(U)$ therefore give
\eqref{tc:mixed} for all $s\in S(U)$.  The local scheduling in
\Cref{seq:diagonal} ensures the same conclusion when a controller is
first selected in a smaller open: dominate it by a finite multiple
of a fixed controller in the chosen larger open, and reduce the
stage error before passing to the next row.  These multiplicities
never change an already fixed cut threshold or lower comparison.

\smallskip\noindent\emph{Recover the square moments.} 
Let $\rho$ now be a regularized limit of coordinate weights $\tau_n$,
with $\sigma=\rho|_Q$ densely finite on the chosen controlling ideal
$Q(U_i)$.  Compact Cuntz support, and the larger controller fixed at
the start, imply that $d_{\tau_n}(b_i)$ is bounded along the defining
ultrafilter.  The constant restrictions converge to $\sigma$ in the
trace topology of \Cref{seq:regularized}.  Bounded-rank continuity
from \Cref{un:trace-cap}, applied also after $F$ using
\eqref{tc:controller}, gives
\[
 \lim_\omega t_i^{(k)}(\tau_n)=t_i^{(k)}(\sigma),\qquad k=1,2.
\]
Consequently \eqref{tc:coordinate-errors} gives convergence of the
uncut square traces to $t_i^{(2)}(\sigma)$, and a bounded first
moment $\lim_\omega\tau_n(h_n(s_i))<\infty$.

There is no loss of square mass on regularization.  For $x\geq0$
and $t>0$, scalar functional calculus gives
\begin{equation}\label{tc:square-mass}
 0\leq\tau(x^2)-\tau((x^2-t)_+)
       \leq\sqrt t\,\tau(x)
\end{equation}
whenever the right side is finite.  Apply this with $x=h_n(s_i)$,
take the ultralimit, and then let $t\downarrow0$.  The definition of
$\rho$ and \eqref{tc:square-mass} yield exactly
\eqref{tc:square}.

\smallskip\noindent\emph{Retain the lower ideal support.} 
Fix one lower datum $(y,\eps,\eta)$ for the test $s=f\otimes a$.
For all sufficiently large $n$, \Cref{tc:rank-average} gives
\[
 d_\tau((y-\eps)_+)
       \leq C\,d_\tau((h_n(s)-\eta)_+)
       \quad\text{for every }\tau\in T(Q),
 \qquad C=2\|y\|^2/\eps^2.
\]
The order determination in \Cref{st:realification} implies
\[
 [(y-\eps)_+]\leq
       \lceil C\rceil[(h_n(s)-\eta)_+].
\]
If $(y-\eps)_+=0$, this comparison has no content and needs no
witness. Otherwise $C>0$; set $N=\lceil C\rceil\geq1$.
The uniform comparison implementation in
\Cref{seq:uniform-witness} gives rectangular columns $w_n$ such that
\[
 \left\|w_n^*\operatorname{diag}_N(h_n(s))w_n-(y-\eps)_+\right\|
       <1/n,
 \qquad
 \|w_n\|^2\leq\frac{\|y\|+1/n}{\eta}.
\]
This bound is independent even of $N$.  The columns pass to the
sequence algebra and put $(y-\eps)_+$ in $\Ideal_D(h(s))$.
Letting the source cut decrease to zero puts $y$ itself in this ideal.
The previously chosen $y$'s and
their cuts generate $Q(U_a)$, proving \eqref{tc:lower-ideal}.
\end{proof}
